Autonomous Surface Vehicles (ASVs) are essential for scalable marine monitoring tasks such as bathymetry, pollution tracking, and border patrol. However, the operational range of these vehicles is strictly limited by battery capacity. High-fidelity perception using sensors like Radar and Lidar is crucial for safe navigation and detailed mapping, but these active sensors deplete energy rapidly. Furthermore, the energy cost of propulsion varies drastically across different water environments; moving against a strong river current consumes far more energy than navigating a calm lake.

This proposal asks a simple question: can an ASV learn to balance its sensor usage and motion strategies to maximize exploration coverage without running out of power? We propose a reinforcement learning (RL) policy that jointly controls vehicle velocity and sensor duty-cycling, explicitly conditioned on the water-body type (`lake`, `river`, `coast`). The core hypothesis is that an agent aware of both its internal energy state and external environmental dynamics will achieve higher information gain per unit of energy than standard energy-blind exploration methods or fixed-schedule sensor baselines.

The method builds on literature in active perception and energy-aware robotics \cite{thrun2005probabilistic,yamauchi1997frontier,charrow2015information,klingensmith2016articulated,sharma2020mineral,viseras2016energy,mccammon2021ocean,liu2022energy,sukkar2019multi,stachniss2005information,chen2019autonomous,lacerda2019probabilistic,vrmac2023energy}. Prior work typically separates path planning from power management or assumes a static environment. Our contribution is a unified policy that treats sensor activation as a learned action and adapts its behavior to specific hydrodynamic contexts (Fig.~
ef{fig:method-overview}), with explicit comparison against related methods (Fig.~
ef{fig:rw-matrix}).
